\section{A common interval for bounds on strict cohomology}\label{sec:cohomology}

The second divisibility estimate requires a weighted sum of strict cohomology lengths over the original auxiliary field. The analytic prime pair has already been fixed in Appendix~\ref{sec:hosts}. We may shorten its interval to control the twists simultaneously while retaining that same pair.
Cyclotomic towers control the lengths outside a bounded
number of exceptional twists, and a height estimate controls
those twists on the smaller interval. Let \(E\) be the auxiliary abelian CM field defined by \eqref{eq:cohom-host-field}. At \(2\) the tower gives a bound over \(E(i)\). Lemma~\ref{lem:dyadic-host} controls restriction from \(E\) to \(E(i)\), yielding a bound on the strict lengths over \(E\) before the twists are summed in Section~\ref{sec:family-window}.

Let \(x\) denote the original cohomology scale \(z\) of
\eqref{eq:cohomwin}. For \(S\) the places above \(\ell\) and the
archimedean places, put \(A_j=\Q_\ell/\Z_\ell(-j)\). The strict
length is
\[
 V(E,\ell,j)=\operatorname{length}_{\Z_\ell}
   \ker\!\left(H^1(G_{E,S},A_j)\longrightarrow
                    \prod_{v\mid\ell}H^1(E_v,A_j)\right).
\]
The fields here are totally complex. Poitou--Tate duality identifies
the length with the length of the corresponding degree-two
localization kernel for \(\Z_\ell(j+1)\).
At the original scale \(x\), take the dyadic index \(e(x)\) so that
\begin{equation}\label{eq:cohom-host-dyadic-level}
 x^{2+8\varepsilon}\le 2^{e(x)}\le2x^{2+8\varepsilon}.
\end{equation}
The auxiliary field \(E\) is cut out by the automorphisms
\(\zeta\mapsto\zeta^a\) with
\begin{equation}\label{eq:cohom-host-field}
 a\equiv\pm1\bmod 2^{e(x)+2},\qquad
 \psi(a)=1,\qquad (a/q)=(a/d)=1.
\end{equation}
Its conductor divides \(\operatorname{lcm}(2^{e(x)+2},N,q,d)\).
It is abelian and CM. In fact \(a=-1\) fails the condition at
\(q\equiv3\bmod8\), so complex conjugation acts nontrivially.

For a fixed interval exponent \(J\) and arbitrarily large
endpoints \(X\), write \([W,W^2]\) for the
subinterval of \([X^{1/J},X]\) on which the bound holds for
every finite set of primes. The proof below establishes
finiteness and the bound in \eqref{eq:cohomwin}. The cutoff
exponent \(401\) leaves room for the constant rescaling
in Section~\ref{sec:family-window}.

\subsection{Localization kernels in cyclotomic towers}

Variation of the cyclotomic layer changes the specialization of one Galois module. Describing the strict kernels as specializations locates the twists for which the usual length bound deteriorates.

Take a base \(L_0\) containing \(\mu_\ell\) for odd \(\ell\) or
\(\mu_4\) for \(\ell=2\). Let \(L_\infty/L_0\) be its cyclotomic
\(\Z_\ell\)-extension and \(L_m\) the layer of degree \(\ell^m\).
Let \(M\) be the abelian pro-\(\ell\) Galois group over \(L_\infty\)
unramified away from \(\ell\) with splitting at all places above \(\ell\).
The coefficient module \(\Q_\ell/\Z_\ell(-j)\) has trivial action
over \(L_\infty\). Denote by \(u\) the cyclotomic value of a topological
generator and by \(A\) its action by conjugation on \(M\).

\begin{lemma}[Lengths away from finitely many cyclotomic twists]\label{lem:tower-length-centres}
Suppose \(M\) is generated by \(R\ge1\) elements over \(\Z_\ell\).
Put \(j_0=\lfloor\log_\ell R\rfloor+3\) and \(c=v_\ell(u-1)\).
There are \(s_1,\ldots,s_R\in\Z_\ell\) such that for every \(m\ge0\),
integer \(j\ge1\), and \(B\ge0\) with
\(\lvert j-s_i\rvert_\ell\ge\ell^{-B}\) for all \(i\),
the strict localization kernel has finite length and
\begin{equation}\label{eq:tower-bound}
 V(L_m,\ell,j)\le R\big(\max(m,j_0)+c+B\big).
\end{equation}
The centres can be chosen once for the whole tower.
\end{lemma}

\begin{proof}
\proofstep{The coinvariant module.}
Multiplication by
\(u^{-j\ell^m}-1\ne0\) is surjective on \(\Q_\ell/\Z_\ell\).
By Lemma~\ref{lem:tower-coinvariants}, the procyclic cohomology computation and
inflation--restriction identify the global degree-one group
with the invariant homomorphisms upstairs. The analogous computation
at the places over \(\ell\) identifies the local splitting condition.
For a continuous homomorphism \(\phi\) upstairs, invariance
under the generator of \(\operatorname{Gal}(L_\infty/L_m)\)
means
\[
 \phi(A^{\ell^m}z)=u^{-j\ell^m}\phi(z)\qquad (z\in M).
\]
Equivalently \(\phi\) kills \(\big((u^j A)^{\ell^m}-1\big)M\).
The strict group is the continuous dual of that quotient.
Taking the Pontryagin dual recovers
\[
   M/\bigl((u^j A)^{\ell^m}-1\bigr)M.
\]
The inverse of the coefficient character occurs when forming
the coinvariant operator.
\proofstep{Growth above a reference layer.}
By continuity the reduction of \(A\) modulo \(\ell\)
is unipotent. Put \(j_0=\lfloor\log_\ell R\rfloor+3\) and
\(g=u^jA\). If \(m\ge j_0\), then \(D=g^{\ell^m}-1\)
maps \(M\) to \(\ell^2M\). The reduction of \(M\) has dimension at most \(R\), so
\[
   g^{\ell^k}-1\equiv(g-1)^{\ell^k}=0\pmod\ell
   \qquad(\ell^k>R).
\]
Another \(\ell\)-th power gives \(g^{\ell^{k+1}}-1\equiv0\pmod{\ell^2}\)
by the binomial formula. This applies for \(m\ge j_0\). Set \(I=DM\).

Since \(D\)
is \(\Z_\ell\)-linear, its powers satisfy
\(D^r M\subset\ell^{2(r-1)}I\) for \(r\ge1\).
Thus \(D'=(1+D)^\ell-1\) takes values in \(\ell I\), and
the terms with \(r\ge2\) take values in \(\ell^2 I\).
Reduction modulo \(\ell^2 I\) gives
\[
      D'M+\ell^2 I=\ell DM+\ell^2I=\ell I.
\]
Apply Nakayama to the finitely generated \(\Z_\ell\)-module \(\ell I/D'M\)
to get \(D'M=\ell I\). The exact sequence
\[
 0\longrightarrow I/\ell I\longrightarrow M/\ell I
       \longrightarrow M/I\longrightarrow 0
\]
then shows that the length of a finite cokernel increases by
at most \(R\), since \(I\), an image of \(M\), is generated
by at most \(R\) elements. Layers below \(j_0\) have cokernels
which are quotients of the reference cokernel.

\proofstep{Distance from a centre at the reference layer.}

Growth along the tower is now bounded in terms of a fixed reference layer. The exceptional twists arise from the eigenvalues of its cyclotomic action, so their distance to the twist under consideration enters here.

Lift \(A^{\ell^{j_0}}\) to a matrix \(A_*\) on a free module on
the \(R\) generators, and put \(w=u^{\ell^{j_0}}\).
The reference cokernel is obtained by the surjection
\[
 \operatorname{coker}(w^j A_*-1)\twoheadrightarrow \operatorname{coker}(w^j A^{\ell^{j_0}}-1\mid M).
\]
 Let \(\alpha_i\) be the eigenvalues of \(A_*\),
counted with multiplicity. Their absolute values are at most one.
If \(\abs{\alpha_i}_\ell=1\), choose \(s_i\in\Z_\ell\) minimizing
\(\abs{w^s\alpha_i-1}_\ell\), which is continuous on the compact
set \(\Z_\ell\). The ultrametric inequality and minimality give
\[
   \abs{w^j\alpha_i-1}_\ell\ge
      \abs{w^j-w^{s_i}}_\ell
      =\ell^{-c-j_0}\abs{j-s_i}_\ell ,
\]
where \(c=v_\ell(u-1)\). The identity follows from the valuation of
integer powers and then by continuity on \(\Z_\ell\).

For \(\abs{\alpha_i}_\ell<1\) the factor equals one and one can
take any \(s_i\in\Z_\ell\).
Smith normal form bounds the length of the reference cokernel by
the valuation of the determinant of \(w^j A_*-1\). Under the
separation hypothesis each of its \(R\) factors has valuation at most
\(c+j_0+B\). For \(m\le j_0\) take the quotient of this cokernel.
For larger \(m\) add at most \(R\) for each further layer. This proves
\eqref{eq:tower-bound}.
\end{proof}

\subsection{Uniform rank at the prime two}

At the prime two the original field need not contain a fourth root of unity. A uniform bound on generators upstairs is useful only after accounting for the kernel of restriction from the original field.

For the prime \(2\) start with the fixed abelian CM field obtained
from the character field by adjoining \(i\).
\begin{lemma}[A uniform number of dyadic generators]\label{lem:dyadic-generators}
Let \(d_0=[K_\psi:\Q]\), and let \(g_0\) bound the \(2\)-ranks of the
class groups in the fixed cyclotomic tower of \(K_\psi(i)\).
For the prime pair in \eqref{eq:cohom-host-field}, after adjoining
\(\sqrt{-q}\) and \(\sqrt{-d}\), the abelian
pro-\(2\) group with the unramified and local splitting conditions
in Lemma~\ref{lem:tower-length-centres} has generator rank \(R\) bounded by 
\[
 R\le4(g_0+6d_0).
\]
 This bound is uniform in the prime pair.
\end{lemma}

\begin{proof}
Ferrero--Washington's theorem gives zero Iwasawa \(\mu\)-invariant
for the cyclotomic \(\Z_2\)-extension of this abelian field
over \(\Q\). The class groups of the finite layers consequently
have bounded \(2\)-rank. The same bound \(g_0\) holds for
unramified quadratic characters of the infinite layer.
For a finite list of such characters, descend to a finite layer
after taking a sufficiently small open cyclotomic subgroup.
One can do this by first adjoining the fields of the characters
over the infinite layer, descending each finite extension to a
finite layer and then enlarging the layer until local inertia
maps trivially at the finitely many ramified places.
If \(V\) is the span of a finite independent list in the infinite layer, the restriction image from that common finite layer contains \(V\). Consequently
\[
 \dim_{\mathbf F_2} V\le
 \dim_{\mathbf F_2}\operatorname{Hom}(\operatorname{Cl}(K'),\mathbf F_2)
 \le g_0 ,
\]
where \(K'\) denotes the finite layer. Taking all finite lists proves the claimed bound.
\hfill\mbox{\cite{ferrero-washington}}

If the fixed character field has degree \(d_0\), the numbers of
primes above \(q,d\) in the fixed tower are at most \(2d_0,4d_0\).
Their decomposition subgroups in the full dyadic tower have
indices
\[
    [\Z_2^\times:\overline{\langle q\rangle}]
      =2^{v_2(q^2-1)-2}=2,\qquad
    [\Z_2^\times:\overline{\langle d\rangle}]
      =2^{v_2(d^2-1)-2}=4 .
\]
Here each projection to the sign quotient is onto, and the square
generates a subgroup of \(1+4\Z_2\) with the displayed index.
Allowing ramification at any one of these primes costs at most
one quadratic character parameter by tame inertia.
The two auxiliary quadratic extensions together have degree at
most four over the fixed tower. Apply the Schreier generator
bound to the subgroup of index at most four in the maximal
pro-\(2\) Galois group with this allowed ramification. Passing to the abelian quotient of this subgroup gives
\[
  R\le 4(g_0+2d_0+4d_0)=4(g_0+6d_0)
\]
in \eqref{eq:tower-bound}.
The splitting condition only decreases this bound. The constant
is thus fixed before the auxiliary prime pair varies.
\end{proof}

For the finite set of odd primes treated by the tower argument,
one takes the rank bound over the fixed abelian base for that
prime.

\begin{lemma}[Restriction at the prime \(2\)]\label{lem:dyadic-host}
Let \(K_\psi\) be the fixed field of \(\ker\psi\) in the
cyclotomic field of its conductor. Let \(L_0\) be the field
\(K_\psi(i,\sqrt{-q},\sqrt{-d})\). For \(E\) defined by
\eqref{eq:cohom-host-field} at dyadic level
\eqref{eq:cohom-host-dyadic-level}, the extension \(E(i)\)
is a layer \(L_m\) of degree \(2^m\) in the cyclotomic
\(\Z_2\)-extension of \(L_0\), with \(m\le e(x)+2\).
For every \(j\ge1\), the localization kernel over \(E(i)\)
controls the kernel over \(E\) by
\begin{equation}\label{eq:dyadic-host}
                  V(E,2,j)\le 1+V(L_m,2,j).
\end{equation}
\end{lemma}
\begin{proof}
\proofstep{Identifying the tower layer.}
Put \(n=e(x)+2\), \(F_n=\Q(\mu_{2^n})\), and write \(F_n^+\)
for its maximal real subfield. In the cyclotomic field containing
the four conductors, the stabilizer of a compositum is the
intersection of the stabilizers of its factors. The signs at \(q,d\)
cut out \(\Q(\sqrt{-q})\) and \(\Q(\sqrt{-d})\), since both primes
are \(3\bmod 4\). The condition \(a=\pm1\bmod 2^n\) cuts out
\(F_n^+\). Thus
\[
 E=K_\psi(\sqrt{-q},\sqrt{-d})F_n^+.
\]
For \(n\ge2\), complex conjugation acts on \(i\in F_n\) by
\(-1\), so \(F_n^+(i)=F_n\).

It follows that \(E(i)=L_0 F_n\).
Adjoining all dyadic roots to \(L_0\supset\mu_4\) gives a
\(\Z_2\)-extension. The subgroup corresponding to \(L_0F_n\)
is therefore the unique subgroup of its index \(2^m\), with
\[
        2^m=[F_n:F_n\cap L_0]\le[F_n:\Q]\le 2^n.
\]
The bound \(m=O(1+\log x)\) is uniform in the prime pairs.
The extension over \(E\) has degree one or two. For the trivial
target, an exponent congruent to \(-1\bmod 2^n\) and
\(1\bmod qd\) fixes \(E\) and moves \(i\), so the quadratic
case already occurs for the zeta application.

\proofstep{Restricting strict classes.}
Set \(A=\Q_2/\Z_2(-j)\). Let \(H^1_{\rm ur(2)}(E,A)\) denote classes unramified
outside the dyadic and infinite places, and use analogous
notation for \(L_m\). Upstairs inertia maps into downstairs
inertia. With \(w\mid v\mid2\), the two localization rows
fit in the commutative diagram
\[
 \begin{array}{ccc}
 H^1_{\rm ur(2)}(E,A)&\xrightarrow{\mathrm{loc}}&
      \displaystyle\prod_{v\mid2} H^1(E_v,A)\\
 \downarrow{\scriptstyle\mathrm{res}}&&\downarrow{\scriptstyle
                    (\mathrm{res}_{w/v})_{w\mid v}}\\
 H^1_{\rm ur(2)}(L_m,A)&\xrightarrow{\mathrm{loc}}&
      \displaystyle\prod_{w\mid2} H^1((L_m)_w,A).
 \end{array}
\]
The left arrow therefore restricts to the strict kernels. View these classes in absolute Galois cohomology
by inflation. For \(H=G_{L_m}\) and \(C=\operatorname{Gal}(L_m/E)\),
inflation--restriction gives the exact sequence
\[
 0\longrightarrow H^1(C,A^H)\xrightarrow{\mathrm{inf}} H^1(E,A)
             \xrightarrow{\mathrm{res}} H^1(L_m,A)^C .
\]
The strict restriction kernel is the intersection of the image
of inflation with the strict group over \(E\). For relative
degree one this kernel is zero.

\proofstep{Bounding the restriction kernel.}

Restriction has already placed the image in the strict kernel upstairs. The downstairs length can thus exceed the upstairs bound only by the length of the kernel of restriction.

In relative degree two put \(B=A^H\) and denote the
involution by \(\sigma\). A cocycle is determined by its value
on \(\sigma\), giving the quotient
\[
    \ker(1+\sigma:B\longrightarrow B)/(\sigma-1)B.
\]
It is killed by two, since \(2b=(1-\sigma)b\) on the numerator.
Every finitely generated subgroup of \(A\) is cyclic of
two-power order. The same property holds for any subquotient,
by lifting a finite set of generators to the subgroup before
taking the quotient. A group with this property killed by two
has at most two elements. The restriction kernel on strict
classes consequently has \(\Z_2\)-length at most one. The image of restriction on strict classes
lies in the strict group upstairs. Additivity of length gives
\eqref{eq:dyadic-host} and finiteness downstairs whenever the
upstairs length is finite. \label{cor:cohomology-cohomology-1}
\end{proof}

\begin{construction}
\begin{tabular}{@{}lll@{}}\toprule
Restriction from \(E\) to \(E(i)=L_m\) & Kernel on strict classes &
Possible length loss\\\midrule
\(i\in E\) & zero & \(0\)\\
\([E(i):E]=2\) & subgroup of \(H^1(C_2,A^{G_{L_m}})\) & \(\le1\)\\\bottomrule
\end{tabular}
\caption{Dyadic descent for a single twist with coefficients \(A=\Q_2/\Z_2(-j)\). Restriction maps strict classes upstairs. Its image has length at most \(V(L_m,2,j)\). The only additional term downstairs comes from its kernel.}\label{fig:dyadic-strict-restriction}
\end{construction}

\subsection{Moving the scale past finitely many twists}
\label{subsec:cohomology-twist-threshold}

The tower estimate can be poor near a finite set of centres. The cutoff in the twist index grows with the real scale, so shortening the scale interval can isolate the exceptions that still require a separate estimate.

\begin{lemma}[Avoiding the exceptional twists on one interval]\label{lem:exceptional-twist-window}
For a finite list of towers as in Lemma~\ref{lem:tower-length-centres},
let \(R_*\) bound the total number of centres. Fix \(0<\eta<401\)
and put
\[
 J=4(2\cdot401/\eta)^{R_*}.
\]
For \(X>1\) set \(Z_*=\lfloor X^{401}\rfloor\). Call an integer
\(1\le j\le Z_*\) exceptional if for at least one centre \(s_i\)
of an \(\ell\)-adic tower it satisfies
\(\lvert j-s_i\rvert_\ell<\ell^{-\lfloor\log_\ell Z_*\rfloor}\).
There exists \(W\ge X^{1/J}\) with \(W^2\le X\) such that, for
every real \(x\in[W,W^2]\), every exceptional \(j<x^{401}\)
satisfies \(j\le x^\eta\). Only \(R_*\) and \(\eta\) enter \(J\).
\end{lemma}

\begin{proof}
Two integers exceptional for the same centre differ by a multiple of
\(\ell^{1+\lfloor\log_\ell Z_*\rfloor}>Z_*\), so they coincide.
There are thus at most \(R_*\) exceptional integers, regardless of
where the centres lie.

Put \(u=\log x\). For an exceptional \(j\), exclude the interval
\[
 [a_j,b_j)=\left[\frac{\log j}{401},\frac{\log j}{\eta}\right).
\]
The inequalities \(x^\eta<j<x^{401}\) imply
\(a_j<u<b_j\). The left endpoint is included in the excluded
interval and the right endpoint is allowed, since there
\(j=x^\eta\). Start with
\(v=(\log X)/J\). If \([v,2v]\) meets \([a_j,b_j)\), replace \(v\)
by \(b_j\). The new \(v\) is larger, so that interval can never
interfere again. Also \(a_j\le 2v\) before the replacement, hence
\(b_j\le(2\cdot401/\eta)v\). After at most \(R_*\) replacements,
\[
 (\log X)/J\le v\le (2\cdot401/\eta)^{R_*}(\log X)/J
                   =(\log X)/4 .
\]
Take \(W=e^v\). The forbidden intervals are open at their right
endpoints: equality \(j=x^\eta\) is permitted.
\end{proof}

For the application take \(0<\varepsilon<1/52\), let
\(C_K\ge1\) be the exponent in the height estimate of
Section~\ref{sec:cohom-height}, and put \(\eta=\varepsilon/(2C_K)\).
Lemma~\ref{lem:dyadic-generators} and the fixed bounds for the
finitely many odd-prime towers give \(R_*\) independently of \(q,d\).
Thus \(J\) is available when the prime pair is chosen. The position
of \(W\) may depend on that pair.

\begin{proposition}[Sum of the dyadic lengths over the original field]\label{prop:dyadic-window-total}
For every sufficiently large selected endpoint \(X\), on the
interval from Lemma~\ref{lem:exceptional-twist-window} one has,
for each real \(x\in[W,W^2]\) and its actual host field \(E\),
\[
  \sum_{1\le j<\lfloor x^{401}\rfloor} V(E,2,j)
                  \le \exp(x^{4\varepsilon/5}).
\]
All lengths in the sum are finite. The lower threshold is uniform in
the selected prime pair at fixed character and \(\varepsilon\).
\end{proposition}

\begin{proof}
\proofstep{Exceptional twists upstairs.}
The degree-two estimate at the reference layer is
\begin{equation}\label{eq:height}
 \log\abs{H^2(G_{L,S},\Z_\ell(j+1))}
      \le \exp\bigl((2j[L:\Q])^{C_K}\log(2\abs{D_L})\bigr)
\end{equation}
for totally complex bases containing \(\mu_\ell\) for odd
\(\ell\), or \(\mu_4\) for \(\ell=2\). Apply it with \(L=L_{j_0}\).
Its degree is uniformly bounded. In proving \eqref{eq:height},
Section~\ref{sec:cohom-height} uses an auxiliary field
\(L^\dagger/L\) whose largest root level \(c^\dagger\) satisfies
\(\ell^{c^\dagger}>2j+1\) and \(r_2(L^\dagger)>2j+1\).
Its relative degree is at most \(([L:\Q]+1)(4j+4)\).
This enlargement is already accounted for in \(C_K\).
The discriminant of \(L_{j_0}\)
has logarithm \(O(1+\log(qd))\) by the conductor--discriminant
formula, and hence \(O(x^{\varepsilon/4})\) on the selected window.

Since the degree is bounded and \(j\le x^\eta\), the new
exponent is
\[
   (2j[L_{j_0}:\Q])^{C_K}\log(2|D_{L_{j_0}}|)
       \ll x^{C_K\eta+\varepsilon/4}
       \ll x^{3\varepsilon/4}.
\]
The last inequalities use the fixed choice of \(\eta\) in the
window.
By Poitou--Tate the degree-two group bounds the strict length
at \(j_0\). Lemma~\ref{lem:tower-length-centres} bounds smaller
layers by a quotient and larger layers by successive growth:
\[
 V(L_m,\ell,j)\le V(L_{j_0},\ell,j)+R\max(0,m-j_0)
       \ll 1+\log x+\exp(Cx^{3\varepsilon/4}).
\]
\proofstep{Restriction and summation.}

The exceptional and nonexceptional twists have been controlled on one upstairs interval. Restriction contributes at most one additional unit of length for each twist before the summation on the original field.

For the other twists \eqref{eq:tower-bound} gives
\(O(1+\log x)\) since \(m=O(1+\log x)\) and
\(\log Z_*\le401J\log x\). Put \(t=\lfloor x^{401}\rfloor\).
Summing \eqref{eq:dyadic-host} over the \(t-1\) integers
\(1\le j<t\) gives
\[
 (\log 2)\sum_{1\le j<\lfloor x^{401}\rfloor} V(E,2,j)
 \le (\log 2)\sum_{1\le j<\lfloor x^{401}\rfloor} V(E(i),2,j)
      +(\log 2)(t-1) .
\]
The additive term is uniform in \(q,d\). The same estimates, after dividing by the fixed positive constant \(\log 2\), bound the logarithm of one plus the unweighted sum by \(O(1+\log x+x^{3\varepsilon/4})\). This is \(o(x^{4\varepsilon/5})\). For large \(x\) the unweighted sum is therefore at most \(\exp(x^{4\varepsilon/5})\), and the weighted contribution is at most \((\log 2)\exp(x^{4\varepsilon/5})\). The same absorption applies to each of the finitely many other towers.
\end{proof}

\subsection{Bounding the cohomology at the reference layer}
\label{sec:cohom-height}

The interval argument leaves a bounded list of twists for which proximity to a cyclotomic centre need not give useful decay. For these twists a bound in terms of the discriminant and the degree of a reference layer is sufficient.

\begin{proposition}[Height bound uniform in the auxiliary layer]\label{prop:uniform-reference-height}
Let \(L\) be the reference layer satisfying the cyclotomic hypotheses of Appendix~\ref{sec:cohomology}, and let \(D_L\) be its discriminant. With \(S\) as in \eqref{eq:height}, the bound 
\[
 \log|H^2(G_{L,S},\mathbb Z_\ell(j+1))|\le\exp\!\left((2j[L:\mathbb Q])^{C_K}\log(2|D_L|)\right)
\]
 is uniform in the further layers used to detect the degree-one localization kernel. The detecting cycle is constructed
in the next subsection.
\end{proposition}

\begin{proof}
\proofstep{Choose a bounded reference layer.}

Write \(u=j\), \(k=2u\) and \(m_0=k+1\).
Enlarge \(L\) first to a reference cyclotomic layer with largest
root level \(c\), such that both \(\ell^c\) and \(r_2(L)\)
exceed \(m_0\). The degree of this enlargement over the
original field is at most \(([L_{\rm orig}:\Q]+1)(4u+4)\).
The layer degrees are powers of \(\ell\). The original field
contains \(\mu_\ell\) for odd \(\ell\) and \(\mu_4\) for \(\ell=2\),
so
\[
              \ell\le[L_{\rm orig}:\Q]+1 .
\]
This bounds the extra jump at the first layer exceeding the two
required thresholds.
The relative discriminant divides a power of the
cyclotomic conductor, and \(\ell^c\le 2[L:\Q]\).
Restriction to the enlarged field is injective on the
degree-one localization kernel, by the earlier injection
into the infinite cyclotomic tower.

\proofstep{Reduce localization to equivariant class-group maps.}

In a further layer \(F=L(\mu_{\ell^\nu})\), write
\(Q_\nu=\ell^\nu\), \(s=[F:L]=\ell^{\nu-c}\),
\(\lambda=1+\ell^c\), and take the generator \(\gamma\) acting
on roots by the power \(\lambda\). Put \(g=\gamma^{-1}\),
\(a=\lambda^u\), and \(\Delta=g-a\) in additive notation.
For a homomorphism \(h\) to the scalar target, equivariance reads
\[
 h(gx)=a h(x)\quad\Longleftrightarrow\quad
 h(\Delta x)=0\qquad (x\in\operatorname{Cl}(F)).
\]
Thus it factors through \(\operatorname{Cl}(F)/\Delta\operatorname{Cl}(F)\). We control their number
by counting classes of weighted norms of elements with
divisor in the image of \(\Delta\).

Let \(A=F^\times/\mu(F)\), written additively, and \(I\) the
free divisor group of \(F\). Set
\[
 W=\operatorname{div}^{-1}(\Delta I)/\Delta A,
 \qquad V=\sum_{i=0}^{s-1}a^i g^{-i}.
\]
As \(g^s=1\), the weighted norm gives
\[
 V\Delta=\sum_{i=0}^{s-1}a^ig^{1-i}
              -\sum_{i=1}^{s}a^ig^{1-i}
          =g-a^sg^{1-s}=(1-a^s)g .
\]
The factors commute, so the same holds for \(\Delta V\).
Since \(a^s\ne1\), this product is injective on each
torsion-free group, and hence so is \(\Delta\). Sending an element of \(W\)
to the class of its unique divisor preimage under \(\Delta\)
gives an exact sequence
\[
 0\longrightarrow U/\Delta U\longrightarrow W
      \longrightarrow\operatorname{Cl}(F)[\Delta]\longrightarrow 0,
 \qquad U=\mathcal O_F^\times/\mu(F).
\]
If the divisor preimage is principal, subtracting \(\Delta\)
of a generator leaves a class of a unit. This gives
exactness at \(W\), and every class killed by \(\Delta\)
lifts by choosing a generator after applying \(\Delta\).

\proofstep{Detect the norm classes by residue evaluations.}

The finite-field detection bound used here is the following.
Put
\[
 b_0=k+m_0(c+m_0),\qquad
 T_i=\left|H_i(\mathrm{GL}_{4u+6}(\mathcal O_L),\Z)_{\mathrm{tors}}\right|.
\]
Among elements of \(F^\times\) with divisor in \(\Delta I\)
the number of distinct classes of \(\ell^{b_0}V\alpha\)
modulo \(Q_\nu\)-th powers is at most
\(Q_\nu^{r_2(L)}T_{k+1}T_k\). The detection uses a
stabilized integral-lattice cycle after reduction modulo
\(Q_\nu\). On a residue field at a prime of relative degree
one it has value
\[
    \left(\prod_{0\le i<m_0,\ i\ne u}
                   (\lambda^u-\lambda^i)\right)v_P(V\alpha),
\]
where \(v_P\) is the quotient character of the residue
multiplicative group modulo \(Q_\nu\)-th powers.
The cycle construction is Proposition~\ref{prop:detecting-cycle-residue}.

Since \(c\ge2\) in the dyadic case, the valuation formula
applies to each \(i\ne u\):
\[
    v_\ell(\lambda^u-\lambda^i)=c+v_\ell(u-i)<c+m_0 .
\]
Write the multiplier, including \(\ell^k\), as \(\ell^e u_0\) with \(u_0\) an \(\ell\)-adic unit. The valuation bounds give \(e\le b_0\). For the difference \(z\) of two residue evaluations,
\[
 \ell^e u_0 z=0\quad\Longrightarrow\quad
 \ell^{b_0}z=\ell^{b_0-e}u_0^{-1}(\ell^e u_0z)=0 .
\]
It forces equality over \(F\) by Chebotarev in the Kummer
extension of a \(Q_\nu\)-th root of the ratio. Indeed a
nonidentity Frobenius there precludes a root in the
residue field, and primes of relative degree greater
than one over \(L\) have density zero among primes of
\(F\). Finitely many primes at which the cycles fail to
reduce do not change this detection.

\proofstep{Bound the homology evaluations.}

Use the standard stabilization \(g\mapsto\operatorname{diag}(g,1)\). The Dedekind domain \(\mathcal O_L\) has Bass stable rank at most two. Indeed, for a unimodular triple \((b,c,d)\) first make \(b'=b+ud\ne0\), unless \(b=d=0\) and \(c\) is already a unit. At each prime dividing \(b'\), either \(d\) vanishes and \(c\) is a unit locally, or one residue for \(v\) must be avoided. The Chinese remainder theorem therefore supplies \(v\) with
\[
      (b+ud,\ c+vd)=\mathcal O_L,\qquad
        \operatorname{sr}(\mathcal O_L)\le2 .
\]

The constant-coefficient case of the homological stability theorem of Randal-Williams and Wahl \cite[Theorem~5.11]{rw-wahl-stability} gives an isomorphism in degree \(i\) whenever \(2i\le a-\operatorname{sr}(\mathcal O_L)-2\). Thus in degree \(i=k+1\) all further stabilizations induce isomorphisms from rank
\[
 2(k+1)+4=4u+6.
\]

The stable homology over \(\Q\) has primitive rank at
most \(r_2(L)\) in the odd degree \(k+1\), by Borel's
calculation for a totally complex field. Use the
block repetition operator with \(\ell\) copies, and for
each \(d=2,\ldots,k+1\) subtract \(\ell^d\) times
ordinary stabilization into the same rank. Let \(R_\ell\) act by block repetition, with comparison in
the common stabilized group. In the stable group the repeated block map factors as the diagonal into \(\ell\) copies followed by block sum.

For a rational primitive class \(z\), each term of the iterated coproduct maps to \(z\) under block sum:
\[
 \Delta^{(\ell)}z=\sum_{i=1}^\ell 1^{\otimes(i-1)}\otimes z\otimes1^{\otimes(\ell-i)},
 \qquad R_\ell z=\ell z.
\] Block repetition respects the product on stable homology because reordering the blocks is conjugation in a common general linear group. Consequently
\[
 R_\ell(z_1\cdots z_d)=\ell^d z_1\cdots z_d
          \qquad(z_i\text{ rational primitive}).
\] Thus
\[
            \prod_{j=2}^{k+1}(R_\ell-\ell^j)\big|_{\text{length }d}=0
              \quad(2\le d\le k+1).
\]
Only the rational primitive part in degree \(k+1\) remains. In degree \(k+1\) the only
remaining rational image thus has dimension at most
\(r_2(L)\).

On additive residue evaluations the operation
multiplies by \(\prod_{d=2}^{k+1}(\ell-\ell^d)\),
which is \(\ell^k\) times a unit. The integral classes
therefore give at most \(Q_\nu^{r_2(L)}T_{k+1}\)
values, by reduction of a group of this rank with
torsion bounded by \(T_{k+1}\). Write \(G_{\rm st}=\mathrm{GL}_{4u+6}(\mathcal O_L)\)
for the group at the stabilized rank. The homology coefficient
sequence in degree \(k+1\) is
\[
 \begin{aligned}
 0\longrightarrow H_{k+1}(G_{\rm st},\Z)/Q_\nu
   &\longrightarrow H_{k+1}(G_{\rm st},\Z/Q_\nu)\\
   &\longrightarrow H_k(G_{\rm st},\Z)[Q_\nu]\longrightarrow0 .
 \end{aligned}
\] The last term has order at most
\(T_k\), so it contributes at most that many cosets of the
image of integral classes. This explains the two torsion
factors in the evaluation count.
\hfill\mbox{\cite{borel}}\label{cor:cohomology-height-1}

Now \(W\to A/Q_\nu A\) given by \(\ell^{b_0}V\) is
well defined because \(\ell^\nu\mid a^s-1\). Its image
has size bounded by the preceding count, which also
bounds the image after quotienting by roots of unity.
If a class is in its kernel, write
\(\ell^{b_0}Vx=Q_\nu y\) and apply \(\Delta\).
The identity \(\Delta V=(1-a^s)g\) gives, after cancellation
in the torsion-free group \(A\),
\[
       \ell^{b_0}\frac{1-a^s}{Q_\nu}\,gx=\Delta y .
\]
As \(v_\ell(a^s-1)=\nu+v_\ell(u)\) and \(g\) is invertible,
the \(\ell\)-primary kernel is killed by \(\ell^{b_0+v_\ell(u)}\).

For \([x]\in W[\ell]\) choose \(y\) with \(\ell x=\Delta y\).
Changing \(x\) by \(\Delta z\) changes \(y\) by \(\ell z\),
by injectivity of \(\Delta\) on \(A\). Thus
\[
 W[\ell]\longrightarrow (A/\ell A)^g,\qquad [x]\longmapsto[y]
\]
is well defined because \(g- a\) kills \(y\bmod\ell\)
and \(a\equiv1\bmod\ell\). If \(y=\ell z\), the torsion-free
identity gives \(x=\Delta z\), proving injectivity.
For \(\operatorname{div}x=\Delta D\), taking divisors of
\(\ell x=\Delta y\) proves
\[
 \Delta(\ell D-\operatorname{div}y)=0
                       \quad\Longrightarrow\quad\operatorname{div}y=\ell D,
\]
since \(\Delta\) is injective on the divisor group.
\proofstep{Control the divisor and unit errors.}

For an invariant class choose a representative with
\(\gamma y/y=z^\ell \rho^t\) in multiplicative notation,
where \(\rho\) generates the roots of unity and
\(0\le t<\ell\). Norms show there are at most
\(\ell^2\) pairs \((t,N_{F/L}z)\). For two elements in one
such part, Hilbert 90 applied to the ratio of their
\(z\)'s shows that their classes differ by a class from
\(L^\times\). This class has valuations divisible by
\(\ell\) away from \(\ell\), because the extension is
unramified there. The classes of such elements modulo
\(\ell\)-th powers are bounded by the order of the
\(\ell\)-torsion in the ideal class group of the
\(\ell\)-integers times the order of their unit quotient
modulo powers. These bounds give
\[
             |W[\ell]|\le\ell^{d_0},
             \qquad d_0=2[L:\Q]+\lfloor\log_2 h_L\rfloor+2 ,
\]
where \(h_L\) is the class number.

\proofstep{Cancel the layer degree in the class group.}

Finally \(g\) acts on the complex places of \(F\) in
\(r_2(L)\) free cycles of length \(s\). The logarithmic
unit lattice spans the hyperplane with coordinate sum
zero, so the determinant of \(g-a\) on it has absolute
value \((a^s-1)^{r_2(L)}/(a-1)\). Its \(\ell\)-valuation is
\[
           (\nu+v_\ell(u))r_2(L)-c-v_\ell(u).
\]
Take \(\ell\)-primary orders in the exact sequence for \(W\).
The units contribute
\(\ell^{(\nu+v_\ell(u))r_2(L)-c-v_\ell(u)}\) to \(|W|_\ell\).
The evaluation count and its kernel bound give
\[
 |W|_\ell\le
 Q_\nu^{r_2(L)}T_kT_{k+1}\ell^{(b_0+v_\ell(u))d_0}.
\]
Dividing by the unit contribution cancels the term
\(\ell^{\nu r_2(L)}\). On the finite class group the
kernel and cokernel of \(\Delta\) have the same order,
also on the \(\ell\)-primary part.

Every equivariant
homomorphism to \(\Z/Q_\nu\Z\) factors through the
cokernel of \(\Delta\), and the number of homomorphisms
from a finite abelian group to \(\Z/Q_\nu\Z\) is at
most the order of its \(\ell\)-primary part. It follows that
the number of these equivariant homomorphisms is bounded by
\[
                    T_kT_{k+1}
                     \ell^{(b_0+v_\ell(u))d_0+c+v_\ell(u)}.
\]
This is independent of \(\nu\). Every finite set of
classes in the localization kernel injects into these
homomorphisms at a sufficiently high common layer.
Their continuous cocycles and the local coboundary
values over \(\ell\) have finite order. In the layer
where they are fixed, localization makes the restricted
homomorphisms everywhere unramified, so they factor
through the class group. The equivariance exponent is
\(a\) for \(g=\gamma^{-1}\). Injection follows upon further
restriction to the infinite tower.
\label{cor:cohomology-height-2}

\proofstep{Bound reference-layer torsion.}

The torsion bound needed for \(T_i\) has the form
\[
                        \log T_i\le
                  (2|D_L|)^{(2u[L:\Q])^C}
                  \qquad(0\le i\le 2u+1)
\]
with absolute \(C\). One obtains it from the lattice
complex of the arithmetic general linear group. Use the
region of positive Hermitian forms in which every nonzero
integral vector has total squared norm at least one.
Its vertices are determined by the active minimum-vector
inequalities. Tuples of vertices with a common minimizing
positive functional give the nerve of the corresponding
normal cones.

The cones cover the open cone of strictly
positive functionals and have convex intersections.
The ordered-tuple complex is acyclic, by the nerve
argument on compact convex subsets of that open cone.
The common minimum vectors must span. Otherwise take a
nonzero positive semidefinite form \(B\) vanishing on their span.
For small \(\epsilon>0\) the active inequalities for \(Q\) remain
equalities for \(Q-\epsilon B\), and the other inequalities remain
valid. The minimizing positive functional \(\lambda\) would then satisfy
\[
           \lambda(Q-\epsilon B)=\lambda(Q)-\epsilon\lambda(B)<\lambda(Q).
\]

For a tuple choose \(n\) independent common minimum
vectors as the columns of a matrix, with \(n=4u+6\) and
\(d=n[L:\Q]\). Their lattice has index at most
\(I=|D_L|^{n/2}d^{d/2}\) by Hadamard's inequality in the
archimedean metrics and Minkowski's lower bound on
the covolume of a lattice of first minimum one.
Hermite normal form counts the row lattices of such
matrices with logarithm of the count \(O(d^2\log(2I))\).
For a vertex form the traces sum to \(n\). The
determinant-product lower bound from the covolume estimate,
together with the bound \(n\) on each eigenvalue, gives
\(\lambda_{\min}\ge I^{-1}n^{-d/2}\) in the chosen frame.

If \(v=(v_\sigma)_\sigma\) is any other minimum vector,
\[
 1=\sum_\sigma Q_\sigma(v_\sigma)
    \ge I^{-1}n^{-d/2}\sum_\sigma \|v_\sigma\|^2 .
\]
Thus \(\|v_\sigma\|\le\sqrt I n^{d/4}\) at each embedding.
For the actual frame index \(I_0\le I\), the coordinates
are in \(I_0^{-1}\mathcal O_L\). For a nonzero \(a\in\mathcal O_L\),
\[
    \max_\sigma |\sigma a|
      \ge \bigl|\operatorname N_{L/\Q}(a)\bigr|^{1/[L:\Q]}\ge1 .
\]
The scaled coordinates therefore have separation at least \(1/I_0\)
in the maximum embedding norm. Packing bounds the number of
possible active vectors. Selecting enough to determine
each vertex bounds the logarithm of the number of
tuple orbits in degree \(i\) by a polynomial in
\(n,d,i\) times \(\log(2|D_L|)\).

The stabilizers are finite and inject into reduction
modulo \(3\), so their orders have the same type of
bound. If \(G\) is the arithmetic group and \(G_\sigma\) a
tuple stabilizer, the equivariant complex yields the page
\[
 E^1_{a,b}=\bigoplus_{\sigma\in G\backslash\mathcal C_a}
                   H_b(G_\sigma,\Z)
 \ \Longrightarrow\ H_{a+b}(G,\Z),
\]
with \(\mathcal C_a\) the ordered tuples of degree \(a\).
Since each stabilizer is finite, over \(\Q\) only \(b=0\)
survives. In the integral degree-zero row a boundary column has at most
\(i+2\) signed faces, counting repetitions, so its Euclidean norm
is at most \(i+2\).

For a minor of size \(a\), Hadamard gives
\[
                         |\det B_a|\le(i+2)^a .
\]
The determinantal divisors bound the torsion of the boundary
cokernels. For a finite stabilizer \(G\) the positive-degree homology is finite and killed by \(|G|\). If such a group \(H_i\) has \(m_i\) generators, then
\[
 (\mathbb Z/|G|\mathbb Z)^{m_i}\twoheadrightarrow H_i,\qquad |H_i|\le |G|^{m_i}.
\]
Bounding their
number of generators by the sizes of the finite bar
groups gives a bound of the same exponential form
for the logarithm of their orders in the required
degrees. The spectral sequence then gives the
displayed bound for \(T_i\).
\label{cor:cohomology-height-3}

The class number in \(d_0\) is bounded by counting integral
ideals of norm at most the Minkowski representative
bound. The enlargement of \(L\) increases its degree and
the logarithm of its discriminant by polynomial factors
in \(u\) and the original degree. Substituting these facts
into the bound for the localization kernel shows that
the logarithm of its order is at most the right side
of \eqref{eq:height}, after increasing \(C_K\).
Write \(K^2\) for the degree-two localization kernel, whose dual is the kernel just bounded by Poitou--Tate duality. The localization sequence gives the order comparison
\[
 |H^2|=|K^2|\,|\operatorname{im}(\mathrm{loc}^2)|
       \le |K^2|\prod_{v\mid\ell}|H^2_v| .
\]
Here the groups are those of the indicated twist. The archimedean terms vanish for the totally imaginary field. Local duality
identifies each order with the order of the fixed
points on the corresponding local torsion twist.
\proofstep{Add the local degree-two terms.}

The global calculation controls the relevant divisor and unit terms. To bound the full degree-two group one must also account for the groups at the places allowed to ramify.


The image of the local cyclotomic character contains
a procyclic subgroup of index bounded by the local
degree. On the twist of exponent \(-j\), its generator
minus one has valuation at most
\(2+\lfloor\log_\ell[L_{\rm orig}:\Q]\rfloor+j\).
Put \(D=[L_{\rm orig}:\Q]\). Each local invariant group is killed by the generator minus one in the rank-one torsion twist. With at most \(D\) local terms, their total contribution to the base-\(\ell\) logarithm is at most
\[
       D(2+\lfloor\log_\ell D\rfloor+j).
\]
Adding this polynomial correction proves
\eqref{eq:height} for the whole degree-two group.
All constants precede the choice of the cyclotomic
layer, the prime and the twist.

\end{proof}

\subsubsection{The detecting cycle}

The height estimate used a class whose residue detects a weighted norm. To justify that detection one must compute the effect of the projector on the valuation at the chosen degree-one prime.

\begin{proposition}[Residue evaluation of the detecting cycle]\label{prop:detecting-cycle-residue}
For \(Q_\nu\), \(\lambda\) and \(V\) as in the reference-layer
argument, the stabilized lattice cycle attached to a weighted norm
has residue value at the degree-one primes used there equal to
\[
 \left(\prod_{\substack{0\le i<m_0\\i\ne u}}
                  (\lambda^u-\lambda^i)\right)v_P(V\alpha).
\]
This is the finite-field evaluation used to compare the
\(Q_\nu\)-power classes.
\end{proposition}

\begin{proof}
\proofstep{The cyclic-group cycle and its weight projector.}

Put \(B=\Z/Q_\nu\Z\) and \(w=\ell^c\). In the bar complex
of the additive cyclic group \(B\) with coefficients
in \(B\), use the degree-\(2u\) cycle
\[
           z=\sum_{b_1,\ldots,b_u\in B}
                      [b_1\mid1\mid\cdots\mid b_u\mid1].
\]
The first face cancels because \(|B|=0\) in the
coefficient ring. For the others, consecutive terms
cancel after translating a summation variable.
For \(0\le t\le m_0\), send it to \(B^2\) by
\(d_t(b)=(b,tb)\). The cyclic periodic resolutions give
\(m_0\) generators in homological degree \(2u\) for
\(B^2\) with these coefficients. Lift the \(m_0+1\)
classes to the free integral module on these generators.
An integral relation whose coefficients have no common
\(\ell\)-factor yields
\[
             \sum_t n_t(d_t)_* z=\partial b
\]
for a bar chain \(b\), with some \(n_{t_0}\) a unit.
\proofstep{Divide the common factor before reduction.}

The division of a common factor is done in the free
integral module before reducing.

Use the group ring of \(B^2\) to separate the linear
functionals \((a,b)\mapsto a+tb\). For this let
\(P(X)=\prod_{i=0}^{m_0-1}(X-\lambda^i)\) and set
\(Q(X)=\sum_i [X^i]P(X)X^{\lambda^i}\).
Put \(J=(X^w-1)\). Substitution \(U_0:X\mapsto X^\lambda\)
is the identity modulo \(J\), and
\[
   U_0(X^w-1)=X^{\lambda w}-1\equiv\lambda(X^w-1)\pmod{J^2}.
\]
It therefore acts by \(\lambda^j\) on \(J^j/J^{j+1}\).
Applying \(P(U_0)\) to \(X\) gives
\(Q(X)\in(X^w-1)^{m_0}\).

In the group ring of \(B\),
each \([t_0-t]-1\) for \(t\ne t_0\) divides
\([w]-1\), as the class of \(w\) lies in the subgroup
generated by \(t_0-t\). Here \(0<|t-t_0|<w\).
Therefore the image of \(Q(X)\) under \(X\mapsto[1]\)
is divisible by the product of those factors.
Multiplying a quotient by
\[
       [-\textstyle\sum_{t\ne t_0}t]\,
            \prod_{t\ne t_0} ([(0,1)]-[(t,0)])
\]
in the group ring of \(B^2\), with the initial bracket in
the first coordinate, gives an element \(\eta\).
For \(\pi_t(a,b)=a+tb\), the \(t\)-indexed factor in the
product maps to \([t]-[t]=0\) when \(t\ne t_0\).
For \(t_0\) the initial translation and the product yield the
factors \([t_0-t]-1\) used to form the quotient. Thus
\[
       (\pi_t)_*\eta =
        \begin{cases}
        0,&t\ne t_0,\\
        \sum_i[X^i]P(X)[\lambda^i],&t=t_0 .
        \end{cases}
\]

\proofstep{Stabilize with lattices carrying the prescribed divisors.}

Let \(\operatorname{div}\alpha=\Delta\mathfrak d\).
On the basis indexed by \(j\in B^2\) of \(F^{B^2}\),
let \(\rho(v)\) act diagonally by the characters
\(\zeta^{j\cdot v}\), with \(\zeta\) a primitive
\(Q_\nu\)-th root. Compare the two stable lattices
consisting respectively of the standard coordinates
and of the fractional ideals with divisors
\(\eta(j)\mathfrak d\).
For \(t\ne t_0\) the total divisor in each
equal-character block for \(\rho d_t\) is zero.
The classification of projective modules over
\(\mathcal O_F\) gives a block isomorphism
of determinant one in each block.

For \(t=t_0\)
move a coordinate with exponent \(\lambda^i\)
to exponent one using \(\gamma^{-i}\). This acts
on its divisor by \(g^i\).
The resulting divisor sum in that block is
\[
             P(g)\mathfrak d=P_u(g)\Delta\mathfrak d .
\]
It is therefore the divisor of
\[
             \beta=P_u(g)\alpha,\qquad
       P_u(X)=\prod_{0\le i<m_0,\ i\ne u}(X-\lambda^i).
\]
In this formula the action is additive on
multiplicative classes. Take an isomorphism with
determinant \(\beta\) in the block of exponent one
and determinant one in the other blocks, then undo
the coordinate changes. In every case the resulting
map \(g_t\) is \(L\)-linear and carries the same pair
of lattices while commuting with \(\rho d_t\).

Conjugate \(\rho\) by \(g_{t_0}^{-1}\) to obtain
\(\rho'\) on the original lattice, and put
\(u_t=g_{t_0}^{-1}g_t\) in its automorphism group.
\proofstep{Correct the conjugations at chain level.}

The bar prism \(\mathcal P_h\) for conjugation by
\(h^{-1}\) satisfies
\(\partial\mathcal P_h+\mathcal P_h\partial
=(\operatorname{conj}_{h^{-1}})_*-1\).
It follows from the relation with boundary \(b\) that
\[
      \rho'_*b-\rho_*b+
                   \sum_t n_t\mathcal P_{u_t}(\rho'd_t)_*z
\]
is a cycle on the original lattice. In the localized
group the composition of the conjugation homotopies
makes it homologous to
\[
 \sum_t n_t\mathcal P_{g_t}(\rho d_t)_*z.
\]

Divide by the unit \(n_{t_0}\) in the coefficient ring
and add a projective complement with trivial action
to put the cycle on a free \(\mathcal O_L\)-lattice.

\proofstep{Evaluate after reduction.}

Reduction converts the projected cycle into the required residue valuation.


For reduction to a finite field \(K\) containing
\(\mu_{Q_\nu}\), take a cyclic logarithm of \(K^\times\)
and its integral carry two-cocycle. Normalize the
logarithm of \(\zeta\) as \((|K^\times|/Q_\nu)\)
modulo \(|K^\times|\). Cup \(u\) copies of the
carry cocycle with the quotient character \(v_P\).
On the prism over exponent \(x\), each summation
over a bar pair telescopes to \(x\) in \(B\).
Interior insertions cancel by symmetry of the
carry cocycle, giving \(x^u v_P(h)\) for a scalar
prism element \(h\).

For a matrix representation transfer along projective
lines by recording the successive scalar changes of
chosen representatives. If
\[
 g e_L=c(g,L)e_{gL}
\]
in the chosen line representatives, then
\[
 c(g_1g_2,L)=c(g_1,g_2L)c(g_2,L).
\]
This identifies each middle face of the transferred bar simplex with the corresponding product face. The two end faces agree after reindexing the lines, so transfer commutes with the differential. Under direct sum, projection to the
lines of the first nonzero block has fibres of
size a power of \(|K|\), which is one in \(B\).
Thus the evaluation adds across blocks and the
trivial complement contributes zero. On a
scalar-character block it is a homomorphism in the
prism element by the prism composition law. Each elementary transvection lies in the commutator subgroup, which this homomorphism kills. The resulting dependence on the matrix \(g\) is therefore through
\[
 g\longmapsto\det g.
\]

This proves the same scalar formula with \(h\)
replaced by its determinant. Only the block with
determinant \(\beta\) contributes. Taking the product
over the conjugate residue factors of \(F\) evaluates
the weighted norm of \(\beta\). More explicitly the
image of \(\zeta\) in the \(i\)-th conjugate factor
is the \(\lambda^i\)-th power of the reference
image, so its scalar cycle contributes \(a^i\)
times the residue evaluation of \(\gamma^i\beta\).
The scalar factors are periodic modulo \(Q_\nu\), since
\[
 a^{i+s}=a^i a^s\equiv a^i\pmod{Q_\nu}.
\]
This permits cyclic reindexing. In the residue evaluation it gives
\[
       v_P(Vg^j\alpha)=a^j v_P(V\alpha).
\] Summing the coefficients
of \(P_u(g)\) therefore gives
\(P_u(\lambda^u)v_P(V\alpha)\), as required in
the detection bound.
\label{cor:cohomology-height-4}
\end{proof}


\subsection{The principal term at odd primes}

For odd residue characteristic the trivial character has to be separated before decomposing into the remaining components. In the indicated parity and congruence range its strict localization kernel over the rationals vanishes.

\begin{lemma}[The principal odd-prime contribution]\label{lem:principal-odd-strict-zero}
For a prime \(\ell\ge5\) and an even integer \(j>0\) with
\(\ell-1\mid2j\), put \(T=\Q_\ell/\Z_\ell(-j)\). For classes unramified away from \(\ell\), strict localization at \(\ell\) has trivial kernel: 
\[
 \ker\!\left(H^1(G_{\mathbb Q,\{\ell,\infty\}},T)\longrightarrow H^1(\mathbb Q_\ell,T)\right)=0.
\]

\end{lemma}

\begin{proof}
Let \(A=\Q_\ell/\Z_\ell(-j)\) and
let \(S_0\) be the subgroup of degree-one cohomology
unramified outside \(\ell\) with zero localization there.
Every class has \(\ell\)-power order because a continuous
cocycle has finite image in the discrete target. Thus
it suffices to show \(S_0[\ell]=0\). For any representative \(c\) of a class killed by \(\ell\),
write \(\ell c=dm\) and take \(m'\in A\) with \(\ell m'=m\).
The cocycle \(c'=c-dm'\) has the same class and
\[
          \ell c'=dm-d(\ell m')=0.
\]
\proofstep{Trivial residual character.}
If \(\ell-1\mid j\), its action on \(A[\ell]\) is
trivial. The cocycle is an abelian homomorphism
unramified outside \(\ell\), hence trivial on the
absolute Galois group of the cyclotomic \(\ell\)-power
tower by Kronecker--Weber. Inflation from the tower
quotient has zero degree-one cohomology on \(A\).
The finite part of the quotient has order prime to \(\ell\).
For a generator \(\gamma\) of the procyclic part the scalar
\(\gamma-1\) on \(A\) is nonzero and surjective. Consequently
\[
 H^1(\overline{\langle\gamma\rangle},A)
 = A/(\gamma-1)A=A/A=0 .
\]
Taking invariants under the finite prime-to-\(\ell\) quotient is exact. The original
class is zero.
\proofstep{Quadratic residual character.}
If \(\ell-1\nmid j\), the residual action is
quadratic. Since \(j\) is even, \(\ell\equiv1\pmod4\)
and it becomes trivial over \(K=\Q(\sqrt\ell)\).
On the decomposition group at \(\ell\), the adjusted
cocycle equals \(dm_\ell\), since its local class
is zero. Choose an inertia element \(s\) acting by \(-1\) on \(A[\ell]\).
Its action minus one is a unit on \(A\), hence
\[
       (s-1)\ell m_\ell=\ell c'(s)=0\quad\Longrightarrow\quad
                  \ell m_\ell=0.
\]
The restriction to \(G_K\) thus vanishes on inertia at \(\ell\),
since \(G_K\) acts trivially on \(A[\ell]\). It gives a
homomorphism from the ideal class group of \(K\)
to \(A[\ell]\). That class group has order less
than \(\ell\).

In fact Minkowski's bound gives an
ideal of norm at most \(\sqrt\ell/2\) in each class,
and the number of ideals of norm \(n\) is at most
the divisor count of \(n\), which is at most \(n\).
For \(B=\lfloor\sqrt\ell/2\rfloor\) this gives
\[
          h_K\le \sum_{n\le B}n=B(B+1)/2<\ell.
\] The homomorphism is therefore zero.
Inflation from \(\operatorname{Gal}(K/\Q)\)
has zero degree-one cohomology since this group
has order prime to \(\ell\). This proves the vanishing of the principal-character contribution for the indicated odd primes and parity.
\label{cor:cohomology-cohomology-2}
\end{proof}

\input{odd-primes}
